% !TEX encoding = UTF-8 Unicode
\documentclass[10pt,a4paper]{article}

% Encoding and Language
\usepackage[T1]{fontenc}
\usepackage[utf8]{inputenc}
\usepackage[ngerman,latin]{babel}

% Page geometry
\usepackage[top=2.5cm,bottom=2.5cm,left=2.5cm,right=2.5cm]{geometry}

% Fonts
\usepackage{lmodern}
\usepackage{microtype}

% Math
\usepackage{amsmath,amssymb,amsfonts}

% Formatting
\usepackage{titlesec}
\usepackage{fancyhdr}
\usepackage{lastpage}
\usepackage{etoolbox}

% Lists and tables
\usepackage{enumitem}
\usepackage{array}
\usepackage{longtable}
\usepackage{booktabs}

% Misc
\usepackage{graphicx}
\usepackage{xspace}
% Custom hangparas environment (replaces hanging package which causes stack overflow)
\newenvironment{hangparas}[2]{\setlength{\hangindent}{#1}\setlength{\hangafter}{#2}}{}

% Remove paragraph indentation, add space between paragraphs
\usepackage{parskip}
\setlength{\parindent}{0pt}
\setlength{\parskip}{0.5em}

% Title and Author
\title{Carl Friedrich Gau\ss\ Werke --- Band~7, Teil~14}
\author{Carl Friedrich Gau\ss}
\date{}

% Header/Footer
\pagestyle{fancy}
\fancyhf{}
\renewcommand{\headrulewidth}{0pt}
\fancyfoot[C]{\thepage}

% Custom commands for old-style formatting
\newcommand{\gesperrt}[1]{\textls[50]{#1}}

% Reduce spacing in table of contents entries
\newcommand{\tocentry}[2]{#1 \dotfill #2}

\begin{document}

\maketitle
\thispagestyle{empty}

\newpage

% ============================================
% Page 641: BEMERKUNGEN. THEORIE DER BEWEGUNG DES MONDES.
% ============================================
\begin{center}
\textbf{BEMERKUNGEN.}\hfill\textbf{THEORIE DER BEWEGUNG DES MONDES.}\hfill\textbf{641}
\end{center}

\vspace{1em}

Im fünften Abschnitt hat Gauss die Differentialgleichung für $p$ nur bis zu den Gliedern vierter Ordnung einschließlich aufgestellt; bei der Entwickelung des Ausdrucks für $\sin 2v' - V_i$ sind indessen einige Glieder dritter Ordnung mitgenommen, weil sie bei der Integration kleine Divisoren erhalten; die beiden Glieder in Klammern sind hier lediglich deswegen hinzugesetzt, weil Gauss bei Aufstellung des darauf folgenden Ausdrucks sie berücksichtigt hat. Der schließliche Ausdruck von $p$ enthält selbstverständlich die Glieder niederer als 5.~Ordnung nicht, welche bei der Integration aus solchen 5.~Ordnung entstehen; einzelne dieser Glieder hatte Gauss zwar schon in seinem Manuscript nachträglich hinzugefügt. Da das aber nur ganz vereinzelt und flüchtig geschehen ist, und auch vielfach noch nicht die richtigen Werthe gegeben waren, so wurden sie nicht mitabgedruckt; auch sonst fanden sich einige Unrichtigkeiten vor, die verbessert worden sind; Gauss hat hier offenbar seine Untersuchungen ziemlich plötzlich abgebrochen, ohne seine Rechnungen bereits geprüft zu haben.

\hfill Brendel.

\vspace{3em}

\begin{flushright}
81
\end{flushright}

\newpage

% ============================================
% Page 642: BERICHTIGUNGEN UND ZUSÄTZE
% ============================================
\begin{center}
\textbf{BERICHTIGUNGEN}\hfill\textbf{UND ZUSÄTZE.}
\end{center}

\vspace{1em}

Seite 59, Zeile 1. statt »imaginäriarm« lies »imaginäriam«.

Seite 362, Zeile 11 statt »Art.~117« lies »Art.~177«.

Seite 547 sind am Schluss des Ausdrucks für $\delta u$ die in den Bemerkungen auf Seite~607 angeführten Glieder hinzuzufügen.

Seite 547, Zeile 18 v.u. statt »$142^{\circ} 23' 36''$« lies »$142^{\circ} 23' 38''$«.

Seite 661, Zeile 3 statt »[6.]« lies »[12.]«.

\newpage

% ============================================
% Page 643: BEMERKUNGEN ZUM SIEBENTEN BANDE
% ============================================
\begin{center}
\textbf{BEMERKUNGEN}\hfill\textbf{ZUM SIEBENTEN BANDE.}
\end{center}

\vspace{1em}

Der vorliegende siebente Band von Gau\ss' Werken bringt zunächst den Neudruck der Theoria motus. Dieselbe ist, bekanntlich bereits durch Schering im Jahre 1871 im Verlage von F.~A.~Perthes in einer den Bänden I---VI der Gesammtausgabe entsprechenden Ausstattung unter Hinzufügung einiger Notizen aus dem Nachlass neu herausgegeben worden; es schien aber trotzdem geboten, einen solchen Abdruck auch in den vorliegenden siebenten Band der Gesammtausgabe aufzunehmen, wobei nach genauer Durchsicht noch einige weitere Incorrectheiten der ersten Ausgabe verbessert werden konnten; auch wurden die numerischen Beispiele einer genauen Nachrechnung unterzogen und eine große Reihe hier neu aufgefundener Fehler verbessert oder, wo dies nicht möglich war, wenigstens angegeben. Ferner wurde das Werk durch eine große Anzahl darauf bezüglicher neuer Notizen aus dem Nachlass bereichert.

Außerdem bringt der Band den bisher noch unveröffentlichten gesammten theoretisch-astronomischen Nachlass von Gau\ss, vor allem Gau\ss' Untersuchungen über die Störungen der Pallas, deren Bekanntmachung von den Fachmännern seit Jahrzehnten mit Spannung erwartet wurde und die auch noch heute großes Interesse bieten werden.

Zwei kleinere Aufsätze, welche bereits in Band~VI aufgenommen wurden, sind wieder mitabgedruckt worden, weil es der Zusammenhang mit den Nachlassnotizen wünschenswerth erschienen ließ, nämlich die Aufsätze über den Zodiakus der Himmelskörper, Seite~313, und über die Tafel zur parabolischen Bewegung, Seite~371.

Die Bearbeitung, wie die Redaction, lag in den Händen von Herrn Brendel in Göttingen.

Bei der Bearbeitung der Ceres- und der Pallasstörungen wurde eine weitgehende Ergänzung der Nachlassnotizen durch den Bearbeiter nothwendig; es sind deshalb die Notizen, für die aus dem Nachlass meist nur einzelne Formeln und Zahlen entnommen werden konnten, in Petit gesetzt. Ebenso wurde des Raumes wegen der Petitsatz für den Druck der größern Zahlentabellen gewählt, auch wenn die Zahlen aus dem Nachlass stammen.

Sonst sind stets, wie in den übrigen Bänden, die Einschaltungen des Bearbeiters durch eckige Klammern [\ ], die Notizen aus dem Briefwechsel u.\,a., die nicht von Gau\ss\ herrühren, in geschweifte Klammern \{~\} gesetzt.

Weitere den Abdruck der einzelnen Notizen betreffende Hinweise findet man in den Bemerkungen des Bearbeiters hinter jeder einzelnen Abtheilung.

\hfill F.~Klein.

\newpage

% ============================================
% Pages 644-650: INHALT
% ============================================
\begin{center}
\Large\textbf{INHALT.}
\end{center}

\vspace{0.5em}

\begin{center}
\textbf{GAUSS WERKE BAND VII.}\\[0.3em]
\textbf{THEORETISCHE ASTRONOMIE.}
\end{center}

\vspace{1em}

\begin{center}
\textbf{\textsc{Theoria Motus Corporum Coelestium in Sectionibus Conicis Solem Ambientium.}}
\end{center}

\vspace{0.5em}

\begin{hangparas}{1em}{1}
\noindent\textsc{Praefatio} \dotfill Seite 3

\vspace{0.3em}

\noindent\textbf{Liber primus.} \textsc{Relationes generales inter quantitates, per quas corporum coelestium motus circa solem definiuntur.}

\noindent\quad Sectio I. Relationes ad locum simplicem in orbita spectantes \dotfill --- 13

\noindent\quad Sectio II. Relationes ad locum simplicem in spatio spectantes \dotfill --- 33

\noindent\quad Sectio III. Relationes inter locos plures in orbita \dotfill --- 110

\noindent\quad Sectio IV. Relationes inter locos plures in spatio \dotfill --- 148

\vspace{0.3em}

\noindent\textbf{Liber secundus.} \textsc{Investigatio orbitarum corporum coelestium ex observationibus geocentricis.}

\noindent\quad Sectio I. Determinatio orbitae e tribus observationibus completis \dotfill --- 155

\noindent\quad Sectio II. Determinatio orbitae e quatuor observationibus, quarum duae tantum completae sunt \dotfill --- 222

\noindent\quad Sectio III. Determinatio orbitae observationibus quotcunque quam proxime satisfacientis \dotfill --- 236

\noindent\quad Sectio IV. De determinatione orbitarum, habita ratione perturbationum \dotfill --- 258

\vspace{0.3em}

\noindent\textsc{Tabulae} \dotfill --- 291

\vspace{0.5em}

\noindent Bemerkungen \dotfill --- 291
\end{hangparas}

\newpage

% Page 645
\begin{center}
\Large\textbf{INHALT.}\hfill\textbf{645}
\end{center}

\vspace{0.5em}

\begin{center}
\textbf{ZUSÄTZE ZUR THEORIA MOTUS CORPORUM COELESTIUM}\\
\textbf{UND ANDERE NACHTRÄGE ZUR ELLIPTISCHEN BEWEGUNG.}
\end{center}

\vspace{0.3em}

\noindent\textbf{Zusätze zur Theoria motus.}

\vspace{0.2em}

\noindent\textit{Nachlass.}

\vspace{0.2em}

\begin{hangparas}{1em}{1}
\noindent Zu Art.~7. Genäherte Formeln für die größte Mittelpunktsgleichung

\noindent Zu Art.~53. Projection einer Planetenbahn auf die Ekliptik

\noindent Zu Art.~70. Zur Berechnung der Parallaxe \hfill Seite 205

\noindent Zu Art.~70---77. Allgemeine Differentialformeln für die geocentrischen Örter der Planeten \hfill --- 296

\noindent Zu Art.~78. Directe Auflösung der Gleichungen $\sin(P-A) = k\sin(Q-B)$, $\sin(P-A') = k'\sin(Q-B')$ \hfill --- 297

\noindent Zu Art.~88---90. Näherungsformel, um aus $\sin \frac{1}{2}\varphi$ den Logarithmen von $\frac{\varphi}{\sin \varphi}$ zu finden; Vorschriften, um den Logar.~Sinus eines kleinen Bogens zu finden \hfill --- 299
\end{hangparas}

\vspace{0.3em}

\noindent\textit{Veröffentlichungen.}

\vspace{0.2em}

\begin{hangparas}{1em}{1}
\noindent Zu Art.~110 und 160. Zur Auflösung der Aufgabe, aus zweien Radiis vectoribus und dem eingeschlossenen Winkel die elliptischen oder hyperbolischen Elemente zu bestimmen \hfill --- 300

\noindent Zu Art.~113 und 177. Druckfehler in Dr.~Gau\ss' \textit{Theoria motus} etc. vom Senator Bar.\ Oriani; über das Integral $\int e^{-t^2} dt = \sqrt{\pi}$ \hfill --- 301
\end{hangparas}

\vspace{0.3em}

\noindent\textit{Nachlass.}

\vspace{0.2em}

\begin{hangparas}{1em}{1}
\noindent Zu Art.~138. Criterium des Falls, wo der die geocentrische Bewegung tangirende größte Kreis durch den Sonnenort geht \hfill --- 303

\noindent Zu Art.~140. Ableitung einiger Gleichungen dieses Artikels \hfill --- 304
\end{hangparas}

\vspace{0.3em}

\noindent\textit{Veröffentlichung.}

\vspace{0.2em}

\begin{hangparas}{1em}{1}
\noindent Zu Art.~182---184. Berichtigungen \hfill --- 307
\end{hangparas}

\vspace{0.3em}

\noindent\textit{Nachlass.}

\vspace{0.2em}

\begin{hangparas}{1em}{1}
\noindent Zu Art.~183. Bestimmung der wahrscheinlichen Fehler der Resultate der Methode der kleinsten Quadrate \hfill --- 307

\noindent Zusammenhang des Initialzustandes mit den Elementen einer Planetenbahn \hfill --- 308

\noindent Über die Verbesserung der Elemente eines Planeten durch vier beobachtete Oppositionen \hfill --- 310

\noindent Differentialformeln bei Berechnung der Oppositionen \hfill --- 312
\end{hangparas}

\newpage

% Page 646
\begin{center}
\Large\textbf{INHALT.}\hfill\textbf{646}
\end{center}

\vspace{0.5em}

\noindent\textbf{II. Über die Zodiaken der Himmelskörper.}

\vspace{0.2em}

\noindent\textit{Veröffentlichung.}

\vspace{0.2em}

\begin{hangparas}{1em}{1}
\noindent Schreiben des Herrn Geheimen Hofraths Gauss an den Herausgeber der Astronomischen Nachrichten \hfill --- 313
\end{hangparas}

\vspace{0.2em}

\noindent\textit{Nachlass.}

\vspace{0.2em}

\begin{hangparas}{1em}{1}
\noindent Zodiakus auf der Himmelskugel \hfill --- 316

\noindent Zodiakus \hfill --- 319
\end{hangparas}

\vspace{0.3em}

\noindent Bemerkungen \hfill --- 322

\vspace{1em}

\begin{center}
\textbf{ZUR PARABOLISCHEN BEWEGUNG.}
\end{center}

\vspace{0.3em}

\noindent\textbf{Zur Cometenbahnbestimmung.}

\vspace{0.2em}

\noindent\textit{Nachlass.}

\vspace{0.2em}

\begin{hangparas}{1em}{1}
\noindent Parabolische Bewegung. Zeit durch zwei Radios vectores und Chorde

\noindent Zur Berechnung der Chorde aus der Zeit und der Summe der beiden Radien

\noindent Zur parabolischen Bewegung. Bestimmung der Elemente durch zwei Radios vectores und Chorde

\noindent Verhältniss von Dreieck und Sector zwischen zwei Radiis vectoribus in der parabolischen Bewegung
\end{hangparas}

\vspace{0.3em}

\noindent\textit{Briefwechsel.}

\vspace{0.2em}

\begin{hangparas}{1em}{1}
\noindent Annäherungsformel für die Zwischenzeit: Gauss an Olbers 1802 Januar 7 \hfill Seite 334

\noindent Gauss an Olbers 1802 Januar 13 \hfill --- 338
\end{hangparas}

\vspace{0.2em}

\noindent\textit{Nachlass.}

\vspace{0.2em}

\begin{hangparas}{1em}{1}
\noindent Allgemeine Theorie der Berechnung der Cometenbahnen, nebst Anwendung auf den zweiten Cometen des Jahrs 1813 \hfill --- 338
\end{hangparas}

\vspace{0.2em}

\noindent\textit{Briefwechsel.}

\vspace{0.2em}

\begin{hangparas}{1em}{1}
\noindent Supplement zur Theoria motus corporum coelestium:

\noindent\quad Gauss an Olbers 1802 Mai 29 \hfill --- 350

\noindent\quad Gauss an Olbers 1802 Juni 14 \hfill --- 350
\end{hangparas}

\newpage

% Continued page 646
\begin{hangparas}{1em}{1}
\noindent\textbf{II. Zur Berechnung der wahren Anomalie in der parabolischen Bewegung.}

\vspace{0.2em}

\noindent\textit{Nachlass.}

\noindent Tabula nova motus parabolici \hfill --- 351

\vspace{0.2em}

\noindent\textit{Briefwechsel.}

\noindent Zur Berechnung der Tafel für die parabolische Bewegung:

\noindent\quad Gauss an Encke 1808 \hfill --- 368

\noindent\quad Gauss an Olbers 1808 Februar 16 \hfill --- 370

\vspace{0.2em}

\noindent\textit{Veröffentlichung.}

\noindent Schreiben des Herrn Hofraths Gauss an den Herausgeber der Astronomischen Nachrichten 1813 April \hfill --- 371
\end{hangparas}

\vspace{0.5em}

\noindent Bemerkungen \hfill --- 378

\vspace{1em}

\begin{center}
\textbf{STÖRUNGEN DER CERES.}
\end{center}

\vspace{0.3em}

\noindent\textbf{I. Erste Methode.}

\vspace{0.2em}

\noindent\textit{Briefwechsel:}

\vspace{0.2em}

\begin{hangparas}{1em}{1}
\noindent Gauss an Olbers 1802 October 12 \hfill --- 377

\noindent Gauss an Olbers 1802 October 24 \hfill --- 378

\noindent Gauss an Gerling 1813 December 28 \hfill --- 378
\end{hangparas}

\vspace{0.2em}

\noindent\textit{Nachlass.}

\vspace{0.2em}

\begin{hangparas}{1em}{1}
\noindent Die Störungen der Planeten \hfill --- 378
\end{hangparas}

\vspace{0.3em}

\noindent\textbf{II. Tafeln für die Störungen der Ceres.}

\vspace{0.2em}

\noindent\textit{Briefwechsel:}

\vspace{0.2em}

\begin{hangparas}{1em}{1}
\noindent Gauss an Olbers 1802 December 3 \hfill --- 396

\noindent Gauss an Olbers 1802 December 21 \hfill --- 396

\noindent Gauss an Olbers 1805 Januar 4 \hfill --- 397

\noindent Gauss an Olbers 1803 März 1 \hfill --- 398

\noindent Gauss an Olbers 1803 April 8 \hfill --- 400

\noindent Gauss an Olbers 1805 Januar 25 \hfill --- 400
\end{hangparas}

\newpage

% Page 647
\begin{center}
\Large\textbf{INHALT.}\hfill\textbf{647}
\end{center}

\vspace{0.5em}

\begin{hangparas}{1em}{1}
\noindent\textbf{III. Zweite Methode.}

\vspace{0.2em}

\noindent\textit{Briefwechsel:} Gauss an Olbers 1805 Mai 10 \hfill Seite 401

\vspace{0.2em}

\noindent\textit{Nachlass:} Ausführung der Rechnung \hfill --- 402

\vspace{0.2em}

\noindent\textit{Briefwechsel:} Gauss an Olbers 1805 Juli 2 \hfill --- 407

\vspace{0.5em}

\noindent Bemerkungen \hfill --- 410
\end{hangparas}

\vspace{1em}

\begin{center}
\textbf{STÖRUNGEN DER PALLAS.}
\end{center}

\vspace{0.3em}

\noindent\textbf{I. Briefwechsel.}

\vspace{0.2em}

\begin{hangparas}{1em}{1}
\noindent Gauss an Olbers 1806 Januar 3 \hfill --- 413

\noindent Gauss an Olbers 1810 August 6 \hfill --- 413

\noindent Gauss an Olbers 1810 October 24 \hfill --- 414

\noindent Gauss an Olbers 1810 November 26 \hfill --- 415

\noindent Gauss an Olbers 1810 November 30 \hfill --- 417

\noindent Gauss an Olbers 1810 December 13 \hfill --- 418

\noindent Olbers an Gauss 1810 December 19 \hfill --- 418

\noindent Gauss an Schumacher 1811 Januar 8 \hfill --- 419

\noindent Olbers an Gauss 1811 Januar 26 \hfill --- 420

\noindent Schumacher an Gauss 1811 Juni 21 \hfill --- 420

\noindent Gauss an Olbers 1811 August 12 \hfill --- 420

\noindent Olbers an Gauss 1812 April 5 \hfill --- 421

\noindent Gauss an Bessel 1812 Mai 5 \hfill --- 421

\noindent Olbers an Gauss 1812 Mai 12 \hfill --- 421

\noindent Gauss an Olbers 1813 April 8 \hfill --- 422

\noindent Gauss an Olbers 1813 Juli 2 \hfill --- 422

\noindent Gauss an Olbers 1814 April 23 \hfill --- 422

\noindent Gauss an Bessel 1814 Mai 18 \hfill --- 424

\noindent Gauss an Olbers 1814 Juni 15 \hfill --- 424

\noindent Gauss an Olbers 1814 September 25 \hfill --- 425

\noindent Gauss an Olbers 1814 December 31 \hfill --- 425

\noindent Olbers an Gauss 1815 Januar 25 \hfill --- 426

\noindent Gauss an Olbers 1816 Januar 8 \hfill --- 426

\noindent Gauss an Olbers 1816 Februar 10 \hfill --- 426

\noindent Olbers an Gauss 1816 März 7 \hfill --- 427

\noindent Gauss an Olbers 1816 Juli 24 \hfill --- 428

\noindent Gauss an Olbers 1817 Februar 15 \hfill --- 428

\noindent Gauss an Olbers 1817 December 2 \hfill --- 429

\noindent Gauss an Encke 1818 März 25 \hfill --- 430

\noindent Gauss an Olbers 1818 März 31 \hfill --- 431

\noindent Gauss an Schumacher 1822 März 17 \hfill --- 431

\noindent Gauss an Gerling 1834 Februar 8 \hfill --- 432

\noindent Encke an Gauss 1834 October 4 \hfill --- 432

\noindent Gauss an Encke 1834 October 12 \hfill --- 433

\noindent Hansen an Gauss 1843 Februar 7 \hfill --- 433

\noindent Gauss an Hansen 1843 März 11 \hfill --- 436

\noindent Gauss an Bessel 1843 März 21 \hfill --- 437
\end{hangparas}

\newpage

% Page 648
\begin{center}
\Large\textbf{INHALT.}\hfill\textbf{648}
\end{center}

\vspace{0.5em}

\begin{hangparas}{1em}{1}
\noindent\textbf{II. Nachlass.} Exposition d'une nouvelle méthode de calculer les perturbations planétaires avec l'application au calcul numérique des perturbations du mouvement de Pallas.

\vspace{0.2em}

\noindent\quad Préface \hfill Seite

\noindent\quad Section I. Mouvement elliptique \hfill ---

\noindent\quad Section II. Variations instantanées des éléments, produites par les perturbations \hfill ---

\noindent\quad Section III. Méthode de calculer les variations finies des éléments pendant un temps limité \hfill ---

\noindent\quad Section IV. Erster Entwurf. Principes de la détermination du mouvement pour un temps illimité \hfill ---

\noindent\quad Section IV. Zweiter Entwurf. Développement des fonctions périodiques en séries \hfill ---
\end{hangparas}

\vspace{0.5em}

\begin{hangparas}{1em}{1}
\noindent\textbf{III. Nachlass.} Specielle Störungen der Pallas durch Jupiter. Erste Rechnung.

\vspace{0.2em}

\noindent\quad Ableitung der Formeln \hfill ---

\noindent\quad Zum Grunde gelegte Elementensysteme \hfill ---

\noindent\quad Berechnung der Störungen \hfill ---

\noindent\quad Bestimmung der Normalelemente \hfill ---
\end{hangparas}

\vspace{0.5em}

\begin{hangparas}{1em}{1}
\noindent\textbf{IV. Nachlass.} Specielle Störungen der Pallas durch Jupiter. Zweite Rechnung.

\vspace{0.2em}

\noindent\quad Berechnung der Störungen \hfill ---

\noindent\quad Bestimmung der Normalelemente \hfill ---
\end{hangparas}

\vspace{0.5em}

\begin{hangparas}{1em}{1}
\noindent\textbf{V. Nachlass.} Allgemeine Störungen der Pallas durch Jupiter. Erste Rechnung.

\vspace{0.2em}

\noindent\quad Formeln für die Variation der Elemente \hfill ---

\noindent\quad\quad Entwickelung einer periodischen Function in eine Reihe \hfill ---

\noindent\quad\quad Praktische Anwendung des Vorigen \hfill ---

\noindent\quad\quad Rechnungsbeispiel \hfill ---

\noindent\quad\quad Entwickelung einer periodischen Function zweier Veränderlicher \hfill ---

\noindent\quad\quad Berechnung und Entwickelung der störenden Kräfte \hfill ---

\noindent\quad\quad Störungen der Neigung und der Länge des Knoten \hfill ---

\noindent\quad\quad Störungen des Perihels \hfill ---

\noindent\quad\quad Störungen des Logarithmen des halben Parameters \hfill ---

\noindent\quad\quad Störungen des Logarithmen der halben großen Axe \hfill ---

\noindent\quad\quad Störungen des ersten Theils der Epoche \hfill ---

\noindent\quad\quad Störungen des zweiten Theils der Epoche \hfill ---

\noindent\quad\quad Bestimmung der Normalelemente \hfill ---
\end{hangparas}

\newpage

% Page 649 (continued)
\begin{center}
\Large\textbf{INHALT.}\hfill\textbf{649}
\end{center}

\vspace{0.5em}

\begin{hangparas}{1em}{1}
\noindent\textbf{VI. Nachlass.} Allgemeine Störungen der Pallas durch Jupiter. Zweite Rechnung.

\vspace{0.2em}

\noindent\quad Berechnung und Entwickelung der störenden Kräfte; Formeln für die Variation der Elemente \hfill ---

\noindent\quad Störungen der Neigung und der Länge des Knoten \hfill ---

\noindent\quad Störungen des Perihels \hfill ---

\noindent\quad Störungen des Excentricitätswinkels \hfill ---

\noindent\quad Störungen des Logarithmen der halben großen Axe \hfill ---

\noindent\quad Störungen des ersten Theils der Epoche \hfill ---

\noindent\quad Störungen des zweiten Theils der Epoche \hfill ---

\noindent\quad Vereinigung der Störungen der Epoche \hfill ---

\noindent\quad Praktische Integration der vom Argument $18\vartheta - 7\eta$ abhängenden Gleichungen \hfill ---

\noindent\quad Rationales Verhältniss der Umlaufszeiten von Pallas und Jupiter \hfill ---

\vspace{0.2em}

\noindent\quad Bestimmung der Normalelemente \hfill Seite 559

\noindent\quad Verbesserung der Jupitermasse \hfill --- 561
\end{hangparas}

\vspace{0.5em}

\begin{hangparas}{1em}{1}
\noindent\textbf{VII. Tafeln für die Störungen der Pallas durch Jupiter.}

\vspace{0.2em}

\noindent\textit{Briefwechsel:} Proberechnung: Gauss an Encke 1816 August \hfill --- 565

\noindent Formeln zur Berechnung der Tafeln \hfill --- 567

\noindent\textit{Briefwechsel:} Schema für die Ausführung der Rechnungen:

\noindent\quad Gauss an Encke 1816 November 7 und 1817 März 3 \hfill --- 568

\noindent\textit{Nachlass:} Tafeln \hfill --- 572
\end{hangparas}

\vspace{0.5em}

\begin{hangparas}{1em}{1}
\noindent\textbf{VIII. Störungen der Pallas durch Saturn, berechnet von Nicolai.}

\vspace{0.2em}

\noindent\textit{Nachlass.}

\noindent Formeln für die Variation der Elemente \hfill --- 578

\noindent\textit{Briefwechsel}

\noindent Störungen des Logarithmen der halben großen Axe, des ersten Theils der Epoche, der Neigung und der Länge des Knoten:

\noindent\quad Nicolai an Gauss 1814 März 14 \hfill --- 580

\noindent Störungen des Excentricitätswinkels und der Länge des Perihels:

\noindent\quad Nicolai an Gauss 1814 am Charfreitage \hfill --- 581

\noindent Störungen des zweiten Theils der Epoche: Nicolai an Gauss 1814 April 22 \hfill --- 583

\noindent Controlle der Rechnungen: Nicolai an Gauss 1814 August 6 \hfill --- 584

\noindent Vereinigung beider Theile der Epoche: Nicolai an Gauss 1815 März 15 und Juli 17 \hfill --- 585
\end{hangparas}

\newpage

% Page 649 (continued) and 650
\begin{center}
\Large\textbf{INHALT.}\hfill\textbf{649}
\end{center}

\vspace{0.5em}

\begin{hangparas}{1em}{1}
\noindent\textbf{IX. Nachlass.} Störungen der Pallas durch Mars.

\vspace{0.2em}

\noindent Formeln für die Variation der Elemente; Berechnung der störenden Kräfte \hfill --- 587

\noindent Entwickelung der störenden Kräfte \hfill --- 589

\noindent Störungen des Logarithmen der halben großen Axe, des ersten Theils der Epoche und des Excentricitätswinkels \hfill --- 593

\noindent Über die Convergenz der Reihen für die störenden Kräfte; andere Methode ihrer Entwickelung durch Einführung eines Factors $\lambda$ \hfill --- 594

\noindent Bestimmung des Factors $\lambda$ \hfill --- 514

\noindent Aufstellung der Reihen für die störenden Kräfte \hfill --- 587

\noindent Andere Methode zur Bestimmung des Factors $\lambda$; Auflösung der Gleichung $a + b\cos 2m = c\cos(M + D)$ \hfill --- 599
\end{hangparas}

\vspace{0.5em}

\noindent Bemerkungen \hfill --- 601

\vspace{1.5em}

\begin{center}
\rule{0.5\textwidth}{0.4pt}
\end{center}

\vspace{1em}

\begin{center}
\textbf{THEORIE DER BEWEGUNG DES MONDES.}
\end{center}

\vspace{0.3em}

\noindent\textit{Nachlass.}

\vspace{0.2em}

\begin{hangparas}{1em}{1}
\noindent Erster Abschnitt. Fundamentalgleichungen für die Bewegung des Mondes

\noindent Zweiter Abschnitt. Integration der Fundamentalgleichungen mit Beiseitesetzung der Störungskräfte der Sonne
\end{hangparas}

\vspace{2em}

\begin{flushright}
VII
\end{flushright}

\newpage

% Page 650
\begin{center}
\Large\textbf{INHALT.}\hfill\textbf{650}
\end{center}

\vspace{0.5em}

\begin{hangparas}{1em}{1}
\noindent Dritter Abschnitt. Erste Annäherung \hfill Seite 620

\noindent Vierter Abschnitt. Zweite Annäherung zur Breite \hfill --- 626

\noindent Fünfter Abschnitt. Zweite Annäherung zur Länge \hfill --- 634
\end{hangparas}

\vspace{0.5em}

\noindent Bemerkungen \hfill --- 641

\vspace{1.5em}

\begin{hangparas}{1em}{1}
\noindent Berichtigungen und Zusätze \hfill --- 642

\noindent Bemerkungen zum siebenten Bande \hfill --- 643
\end{hangparas}

\vspace{3em}

\begin{center}
Göttingen, Druck der Dieterichschen Univ.-Buchdruckerei (W.~Fr.~Kaestner).
\end{center}

\end{document}
